% Standalone section fragment. Uses amsmath, amsthm, amssymb only.
\section{Bounded selection of regular heights}
\label{sec:regular-heights}

A sentence such as ``move the horizontal edge a little until it misses all
zeros'' is not an algorithm with a primitive recursive bound.  This section
supplies that bound.  All the constants are intentionally generous.
The polynomial facts used below concern rational complex polynomials; their
algebraic roots are represented by finite real-algebraic data.  In particular,
root multiplicities and comparisons with a rational circle are decidable by
bounded algebraic algorithms.  The uniform, variable-degree version of this
finite polynomial lemma is required here, not merely a separate theorem for
each fixed degree.

\subsection{A finite polynomial substitute for Jensen's inequality}

Let $P$ be a polynomial, $P(0)\ne0$, and suppose $|P(z)|\le B$ on
$|z|=R$.  If $a_1,\ldots,a_q$ are any of its roots in $|z|\le r<R$,
repeated according to multiplicity, put
\[
 B_a(z)=\frac{R(z-a)}{R^2-\overline a z},\qquad
 H(z)=P(z)\prod_{j=1}^q B_{a_j}(z)^{-1}.
\]
The apparent divisions cancel polynomial factors, so $H$ is a polynomial.
On $|z|=R$ each $|B_a(z)|=1$.  The constant-coefficient mean identity for a
polynomial therefore gives
\[
 |P(0)|\prod_{j=1}^q\frac R{|a_j|}=|H(0)|\le B,
 \qquad
 |P(0)|(R/r)^q\le B.                         \tag{J}
\]
There is no assumption that the outer circle avoids roots.  The mean identity
is verified by integrating the finitely many monomials, or by a sufficiently
fine finite Fourier sum.  Thus (J), unlike a qualitative invocation of
Jensen's formula, is already a finite algebraic inequality with explicit
approximation margins.

\subsection{A minimum-modulus estimate independent of the degree}

Suppose $P\in\mathbb Q(i)[z]$ satisfies
\[
 |P(0)|\ge\tfrac18,\qquad |P(z)|\le2M\quad(|z|\le4),\qquad M\ge1.
\]
Let $k$ be the least integer for which $2^k\ge16M$, and put $Q=4k$.
By (J), the number $q$ of roots in $|z|\le3$ is at most $Q$:
indeed $(4/3)^q\le16M\le2^k$, whereas $(4/3)^4>2$.
Remove exactly these roots with the radius-four factors $B_a$ above.
The resulting polynomial $H$ has no root in $|z|<3$, satisfies
$|H|\le2M$ on $|z|\le4$, and $|H(0)|\ge1/8$.

For completeness, the only harmonic estimate needed is the following
explicit version of Harnack's inequality.  On $|z|\le11/4$ the function
\[
 u(z)=\log\frac{2M}{|H(z)|}
\]
is nonnegative and harmonic.  The Poisson kernel at $|z|\le5/2$, relative to
the circle of radius $11/4$, lies between positive bounds and is at most
\[
 \frac{11/4+5/2}{11/4-5/2}=21.
\]
Consequently $u(z)\le21u(0)$ and
\[
 |H(z)|\ge \frac{1}{2^{83}M^{20}}\quad(|z|\le5/2).       \tag{H}
\]
This use of Poisson integration has a finite implementation.  If $d=\deg H$,
all roots of $H$ have modulus greater than three, whence its factorization
already gives
\[
 |H(z)|\ge |H(0)|\,12^{-d}\quad (|z|\le11/4).
\]
That inferior but explicit lower bound controls the quadrature and logarithm
errors while proving the much stronger bound (H).  More concretely, write
$H(z)=H(0)\prod_{\nu=1}^d(1-z/b_\nu)$, where $|b_\nu|>3$.
On the radius-$11/4$ circle the logarithm series
$\log(1-z/b_\nu)=-\sum_{j\ge1}(z/b_\nu)^j/j$ has ratio at most $11/12$.
Truncation after $K$ terms has total error at most
$12d(11/12)^{K+1}$, and the angular derivative of $\log|H|$ is at most
$11d$.  Finite Fourier identities prove the mean and Poisson identities
for the truncated sums; the displayed geometric tail and the positive
lower bound permit passage to any requested rational error.
Their dependence on $d$ is harmless: the approximation degree has already
been chosen.  They do not enter the approximation tolerance used to choose
that degree.

If $|z-a_j|\ge\delta$ for every removed root and $|z|\le5/2$, then
\[
 |B_{a_j}(z)|\ge\frac{4\delta}{16+3(5/2)}>\frac\delta6.
\]
For $0<\delta\le1$ this proves
\[
 |P(z)|\ge \frac{1}{2^{83}M^{20}}(\delta/6)^Q
                                                        \tag{MM}
\]
outside at most $Q$ disks of radius $\delta$.  The exponent is controlled by
$\log M$, not by the approximation degree.  This distinction eliminates
an otherwise circular choice of precision.

\subsection{The regular-height algorithm}

Take rational $T\ge8$, write $A=\lceil T\rceil$, and define
\[
 \begin{aligned}
 c&=2+i(T+\tfrac12),& M&=4(A+10)^4,\\
 k&=\min\{j:2^j\ge16M\},& Q&=4k,\\
 \delta&=\frac1{3(Q+2)},&
 L&=\frac1{2^{83}M^{20}[18(Q+2)]^Q}.
 \end{aligned}                                             \tag{RH1}
\]
Every operation in (RH1), including the bounded search for $k$, is primitive
recursive.  Euler--Maclaurin with two correction terms gives
\[
 |\zeta(c+z)|\le M\quad(|z|\le4),\qquad
 |\zeta(c)|\ge\Re\zeta(c)\ge\tfrac14.                    \tag{RH2}
\]
Here the second inequality follows directly from
$\sum_{n\ge2}n^{-2}\le1/4+\int_2^\infty x^{-2}dx=3/4$.
For the first, $\Re(c+z)\ge-2$, $|c+z+j|\le A+10$ for $0\le j\le3$,
and the fourth Bernoulli remainder is bounded by
$|(c+z)_4|/[24(\Re(c+z)+3)]$ after using $|\widetilde B_4|\le1$.
The pole is outside the disk and its distance is bounded below.
These estimates are much smaller than the displayed value of $M$.

Construct a rational complex polynomial $P$ such that
\[
 \sup_{|z|\le4}|P(z)-\zeta(c+z)|\le L/4.                 \tag{RH3}
\]
This is a bounded construction.  One possible recipe uses the Taylor series
at $c$ and the radius-five bound $M_5=8(A+14)^6$, obtained by
Euler--Maclaurin with three correction terms.  Choose $p,b$ with
$2^{-p}\le L/4$ and $M_5\le2^b$, and take
\[
 D=5(p+b+4).
\]
Since $(4/5)^5<1/2$, the Taylor tail beyond degree $D$, on radius four,
is less than $2^{-p}/2$.  Approximating each coefficient to
$2^{-p}/[2(D+1)4^D]$ contributes at most the other half.
Cauchy coefficient integrals can themselves be computed by bounded
quadrature on radius five; the radius-six bound
$M_6=32(A+16)^6$ supplies derivative bounds.  For the $j$th coefficient the
integrand, parametrized on $[0,1]$, has derivative at most
$8(5M_6+jM_5)$ before division by $5^j$.  These explicit bounds give a
primitive recursive sample count.  Thus (RH3) does not invoke an unbounded
search for a successful approximant.

Apply (MM) to $P$.  The possible bad disks have radius $\delta$, and there
are at most $Q$ of them.  Consider the $Q+1$ rational heights
\[
 U_j=T+\frac{j+1}{Q+2},\qquad 0\le j\le Q.              \tag{RH4}
\]
Their spacing is $1/(Q+2)>2\delta$.  Each bad disk can meet at most one of
the horizontal segments $[1/2,2]+iU_j$.  The finite pigeonhole principle
therefore supplies a segment meeting no bad disk.  That entire segment lies
within distance $5/2$ of $c$, so (RH3) and (MM) imply
\[
 |\zeta(\sigma+iU_j)|\ge3L/4\quad(1/2\le\sigma\le2).
                                                               \tag{RH5}
\]

The algorithm need not compute the algebraic roots of $P$.  It simply tests
all the heights (RH4), using a fixed finite mesh.  Let
$J=\lceil12M/L\rceil$ and sample
$\sigma_\ell=1/2+3\ell/(2J)$, $0\le\ell\le J$.
At every sample compute a rational complex value $w_\ell$ with error at most
$L/16$.  Accept a height if
\[
 |w_\ell|^2\ge(5L/8)^2\quad\hbox{for every }\ell.        \tag{RH6}
\]
This is a decidable rational test.  A height satisfying (RH5) passes (RH6).
Conversely, Cauchy's estimate gives $|\zeta'|\le M$ on every candidate
segment, and the mesh spacing is at most $L/(8M)$.  An accepted height has
\[
 |\zeta(\sigma+iU)|\ge7L/16>L/4
 \quad(1/2\le\sigma\le2).                               \tag{RH7}
\]
Return the first accepted height, its mesh, and the margin $L/4$.
There are exactly $Q+1$ candidates and $J+1$ samples per candidate; all
precision and summation bounds have been specified in advance.
This proves a uniform primitive recursive regular-height construction.

\subsection{An explicit certificate for the complete rectangle}
\label{sec:complete-contour}

The accepted horizontal segment is enough to construct every part of the
zero-counting contour with a predetermined precision.  Let its height be
$U\ge8$, and let $L\le1$ be the rational number from the regular-height
construction, so that
$|\zeta(\sigma+iU)|\ge L/4$ for $1/2\le\sigma\le2$.
Put
\[
 \begin{aligned}
 D&=\lceil U\rceil+2,
 & M_\Gamma&=D\,4^{D+D^2},\\
 R&=\lceil U\rceil+3,
 & B&=2(2R+1)R!+R(R+1)(R+1)!,\\
 \mu&=\frac{L}{32M_\Gamma}.
 \end{aligned}                                             \tag{FC1}
\]
The reciprocal-gamma estimate from the analytic foundation gives
$|\Gamma(z)|\ge M_\Gamma^{-1}$ whenever
$\Re z>0$ and $|z|\le D$.
The derivative estimate for the entire completed function gives
$|\xi'(s)|\le B$ on $|s|\le R$.

On the top half-edge $s=\sigma+iU$, $1/2\le\sigma\le2$, we have
$|s(s-1)|\ge1$, $\pi^{-\sigma/2}\ge1/4$, and $|s/2|\le D$.
The defining product therefore gives
\[
 |\xi(s)|\ge \tfrac12\cdot1\cdot\tfrac14
                 \cdot M_\Gamma^{-1}\cdot\tfrac L4=\mu.
\]
On the right edge $s=2+it$, $0\le t\le U$, one has
$|s(s-1)|\ge2$ and $|\zeta(s)|\ge\Re\zeta(s)\ge1/4$.
Consequently $|\xi(s)|\ge1/(16M_\Gamma)\ge\mu$ there.
On the bottom edge $\xi(s)\ge3/8\ge\mu$.
Finally the identity
$\xi(1-\overline s)=\overline{\xi(s)}$
reflects the right edge to the left and the top half-edge to the remaining
top half.  It follows that the complete rectangle has the uniform rational
margin
\[
 |\xi(s)|\ge\mu
 \quad\bigl(s\in\partial([-1,2]+i[0,U])\bigr).          \tag{FC2}
\]
The gamma estimate was used only where $\Re(s/2)>0$; no division by a gamma
value at a pole is hidden in the reflected half.

Divide each horizontal edge into
\[
 J_x=\left\lceil\frac{24B}{\mu}\right\rceil
\]
subintervals and each vertical edge into
\[
 J_y=\left\lceil\frac{8UB}{\mu}\right\rceil
\]
subintervals.  Each segment length is at most $h=\mu/(8B)$.
At every contour vertex $s_j$, evaluate $\xi(s_j)$ by the entire theta
representation to a rational complex number $v_j$ with error at most
$\mu/16$.  The required binary precision is the least $p$ with
$2^{-p}\le\mu/16$; the analytic approximation routine supplies its already
proved primitive recursive tail and quadrature bounds.  Use the same
approximation at the repeated initial/final vertex.

For a contour segment parametrized by
$s(t)=(1-t)s_j+ts_{j+1}$, let
$v(t)=(1-t)v_j+tv_{j+1}$.
The derivative bound implies
\[
 \begin{aligned}
 |\xi(s(t))-v(t)|
 &\le (1-t)Bt|s_{j+1}-s_j|
       +tB(1-t)|s_{j+1}-s_j|+\mu/16\\
 &\le Bh/2+\mu/16\le\mu/8.
 \end{aligned}
\]
The straight homotopy between the analytic image and this rational polygon
therefore has modulus at least $7\mu/8$ everywhere.
The polygon has exactly $2(J_x+J_y)$ segments; its winding integer is
computed by finite rational arithmetic and equals the analytic contour
count.

For example, no trigonometric or branch computation is needed to find that
integer.  For each oriented edge from $a$ to $b$, add one if
$\Im a\le0<\Im b$ and $\det(a,b)>0$; subtract one if
$\Im b\le0<\Im a$ and $\det(a,b)<0$; add zero otherwise.
Here $\det(a,b)=\Re a\Im b-\Im a\Re b$.
The half-open conventions handle vertices on the real axis.
If an edge straddles that axis and its determinant vanishes, it would pass
through the origin, which (FC2) and the homotopy bound exclude.
Thus all branches are decidable exact rational comparisons.
This completes the primitive recursive contour constructor, rather than
merely assuming the existence of the other three zero-free edges.

\section{The zero-counting argument}
\label{sec:counting}

Let
\[
 \xi(s)=\tfrac12s(s-1)\pi^{-s/2}\Gamma(s/2)\zeta(s).
\]
The analytic foundation supplies its entire continuation and the identities
$\xi(1-s)=\xi(s)$ and $\xi(\overline s)=\overline{\xi(s)}$.
For a rational regular height $U\ge8$ equipped with a positive margin on
$[1/2,2]+iU$, let $N(U)$ denote the winding count of $\xi$ on the
positively oriented rectangle $[-1,2]+i[0,U]$.  Its classical interpretation
is the number of nontrivial zeros with $0<\Im\rho<U$, counted with
multiplicity.  The supplied margin, the functional equation, and the
explicit nonvanishing bounds on the other edges turn this contour into a
finite rational winding certificate.

\subsection{The exact argument identity}

The function $\xi$ is positive on the real interval $[-1,2]$.
An explicit lower bound $1/4$ follows from the theta integral
\[
 \xi(s)=\frac12+\frac{s(s-1)}2
 \int_1^\infty \Bigl(\sum_{n\ge1}e^{-\pi n^2x}\Bigr)
       (x^{s/2}+x^{(1-s)/2})\,\frac{dx}{x}.
\]
For $s\in[0,1]$ the integral is at most one: use
$\sum_{n\ge1}e^{-\pi n^2x}\le2e^{-3x}$ and
$x^{s/2}+x^{(1-s)/2}\le2x$.
Thus $\xi(s)\ge3/8$ there; outside $[0,1]$ the integral term is
nonnegative.

Start the continuous argument at zero at $s=2$, and follow the path
$2\to2+iU\to1/2+iU$.  Reflection
$s\mapsto1-\overline s$ and conjugation show that the change of argument
on the whole rectangle is twice the change on this path.  The latter is
\[
 \pi+\vartheta(U)+\arg\zeta(1/2+iU),
\]
where
\[
 \vartheta(U)=\Im\log\Gamma(1/4+iU/2)-\tfrac U2\log\pi
\]
uses the logarithm of gamma continuous in the right half-plane, real on the
positive real axis.  In particular, it is not a repeatedly reset principal
argument.  The factor $s(s-1)$ contributes exactly $\pi$: at the endpoint
its value is $-U^2-1/4$, and both individual arguments were continued through
the upper half-plane.  Therefore
\[
 N(U)=1+\frac{\vartheta(U)}\pi+S(U),\qquad
 S(U)=\frac1\pi\arg\zeta(1/2+iU).                       \tag{C1}
\]
All argument identities can equivalently be read as identities of the
integer winding numbers of rational polygonal approximations, whose
homotopies have the supplied positive margins.

\subsection{Backlund's estimate with a finite polynomial crossing count}

Euler summation gives, for $\Re s>0$,
\[
 \zeta(s)=\frac{s}{s-1}-s\int_1^\infty\{x\}x^{-s-1}\,dx.
\]
Consequently, if $U\ge8$ and $|z-2|\le7/4$, then
\[
 |\zeta(z\pm iU)|\le5(U+4).
\]
Indeed $\Re z\ge1/4$, $|z\pm iU|\le U+4$, and
$|z\pm iU-1|\ge U-7/4>2$.
At $z=2$, the real part of $\zeta(2+iU)$ is at least $1/4$.

Approximate $w\mapsto\zeta(2+w+iU)$ by a rational polynomial $P(w)$
on $|w|\le7/4$, with error less than both $1/8$ and one quarter of the
supplied horizontal nonvanishing margin.  Let
\[
 Q(w)=\tfrac12\bigl(P(w)+\overline{P(\overline w)}\bigr).
\]
This is a real polynomial, $Q(0)\ge1/8$, and
$|Q(w)|\le6(U+4)$ on the outer disk.  By (J), the number $q$ of its
roots in $|w|\le3/2$, counted with multiplicity, satisfies
\[
 q\log(7/6)\le\log(48(U+4)),\qquad
 q\le7\log(48(U+4)),                                   \tag{C2}
\]
because $\log(7/6)\ge1/7$.
The real roots of $Q$ partition the segment $[-3/2,0]$ into at most
$q+1$ intervals.  On each interval $P$ stays in one open half-plane, so its
continuous argument changes by less than $\pi$.
Homotopy from $P$ to $\zeta$ changes only the endpoint arguments, by a
total of less than $\pi$; the vertical path at real part two stays in the
right half-plane.  Allowing an extra unit for all endpoint conventions gives
\[
 |S(U)|\le7\log(48(U+4))+4.                             \tag{C3}
\]
This proof uses only the finitely many roots of a rational real polynomial,
not a procedure deciding the zeros of the real part of an analytic
function.  Multiple roots cause no problem and only enlarge the estimate.

\subsection{Riemann--von Mangoldt with an explicit error}

The explicit complex Stirling estimate proved in the analytic foundation is
\[
 \left|\vartheta(U)-\left(\frac U2\log\frac U{2\pi}
                    -\frac U2-\frac\pi8\right)\right|
 \le1\qquad(U\ge8).
\]
Combining this with (C1)--(C3), and using $\log48<4$, gives the deliberately
coarse but uniform statement
\[
 \boxed{\quad
 \left|N(U)-\frac U{2\pi}\log\frac U{2\pi}
                 +\frac U{2\pi}\right|
 \le8\log(U+4)+40.
 \quad}                                                 \tag{RVM}
\]
The quantifiers are important: $U$ is supplied with a finite nonvanishing
certificate, or is returned by the bounded regular-height algorithm.
No exact decision whether an arbitrary real height is a zero ordinate is
used.

For an integer $n\ge8$, let $U_n\in(n,n+1)$ be the deterministically
selected height and let $C(n)=N(U_n)$.  The contour construction and winding
computation make $C$ a total primitive recursive integer function.
Moreover $C(n+1)\ge C(n)$ and, with
$\mathcal M(t)=t\log(t/(2\pi))/(2\pi)-t/(2\pi)$,
\[
 |C(n)-\mathcal M(n)|\le9\log(n+5)+40.                           \tag{RVM-PR}
\]
This is one particularly simple arithmetical formulation of the classical
asymptotic.  In the classical interpretation, zero counts at any intermediate
height are bracketed by the adjacent certified counts.  A finite-precision
constructive statement uses those brackets rather than a discontinuous
real-to-integer counting function.

